% !TeX program = pdflatex
\documentclass[11pt,a4paper]{amsart}

\newcommand{\notelanguage}{finnish}
% Shared preamble for course notes and exercise sets.

\usepackage[T1]{fontenc}
\usepackage{lmodern}
\providecommand{\notelanguage}{english}
\usepackage[finnish,english]{babel}
\AtBeginDocument{\expandafter\selectlanguage\expandafter{\notelanguage}}

\newcommand{\notesfinnishlanguage}{finnish}
\ifx\notelanguage\notesfinnishlanguage
\newcommand{\notestheoremname}{Lause}
\newcommand{\noteslemmaname}{Lemma}
\newcommand{\notespropositionname}{Propositio}
\newcommand{\notescorollaryname}{Seurauslause}
\newcommand{\notesdefinitionname}{Määritelmä}
\newcommand{\notesexamplename}{Esimerkki}
\newcommand{\notesproblemname}{Tehtävä}
\newcommand{\notesexercisename}{Tehtävä}
\newcommand{\notesremarkname}{Huomautus}
\newcommand{\notessolutionname}{Ratkaisu}
\else
\newcommand{\notestheoremname}{Theorem}
\newcommand{\noteslemmaname}{Lemma}
\newcommand{\notespropositionname}{Proposition}
\newcommand{\notescorollaryname}{Corollary}
\newcommand{\notesdefinitionname}{Definition}
\newcommand{\notesexamplename}{Example}
\newcommand{\notesproblemname}{Problem}
\newcommand{\notesexercisename}{Exercise}
\newcommand{\notesremarkname}{Remark}
\newcommand{\notessolutionname}{Solution}
\fi

\usepackage[a4paper,margin=30mm]{geometry}
\usepackage{microtype}
\usepackage{graphicx}
\usepackage{enumitem}

\newcommand{\exerciseheading}[2]{%
  \par\noindent
  {\large\bfseries #1\par}
  \vspace{0.1em}
  \noindent{\large #2\par}
  \vspace{1.5em}
}

\usepackage{amsmath}
\usepackage{amssymb}
\usepackage{amsthm}
\usepackage{mathtools}
\usepackage{aliascnt}

\usepackage[hidelinks,pdfusetitle]{hyperref}

\numberwithin{equation}{section}
\setlist{itemsep=0.25em,topsep=0.5em}

\theoremstyle{plain}
\newtheorem{theorem}{\notestheoremname}[section]
\newaliascnt{lemma}{theorem}
\newtheorem{lemma}[lemma]{\noteslemmaname}
\aliascntresetthe{lemma}
\newaliascnt{proposition}{theorem}
\newtheorem{proposition}[proposition]{\notespropositionname}
\aliascntresetthe{proposition}
\newaliascnt{corollary}{theorem}
\newtheorem{corollary}[corollary]{\notescorollaryname}
\aliascntresetthe{corollary}

\theoremstyle{definition}
\newaliascnt{definition}{theorem}
\newtheorem{definition}[definition]{\notesdefinitionname}
\aliascntresetthe{definition}
\newaliascnt{example}{theorem}
\newtheorem{example}[example]{\notesexamplename}
\aliascntresetthe{example}
\newaliascnt{problem}{theorem}
\newtheorem{problem}[problem]{\notesproblemname}
\aliascntresetthe{problem}
\newtheorem{exercise}{\notesexercisename}

\theoremstyle{remark}
\newaliascnt{remark}{theorem}
\newtheorem{remark}[remark]{\notesremarkname}
\aliascntresetthe{remark}

\newenvironment{solution}
{
\begin{proof}[\notessolutionname]}
  {
\end{proof}}

\usepackage[nameinlink,noabbrev]{cleveref}

\crefname{theorem}{theorem}{theorems}
\crefname{lemma}{lemma}{lemmas}
\crefname{proposition}{proposition}{propositions}
\crefname{corollary}{corollary}{corollaries}
\crefname{definition}{definition}{definitions}
\crefname{example}{example}{examples}
\crefname{problem}{problem}{problems}
\crefname{exercise}{exercise}{exercises}
\crefname{remark}{remark}{remarks}

\ifx\notelanguage\notesfinnishlanguage
\crefname{theorem}{lause}{lauseet}
\crefname{lemma}{lemma}{lemmat}
\crefname{proposition}{propositio}{propositiot}
\crefname{corollary}{seurauslause}{seurauslauseet}
\crefname{definition}{määritelmä}{määritelmät}
\crefname{example}{esimerkki}{esimerkit}
\crefname{problem}{tehtävä}{tehtävät}
\crefname{exercise}{tehtävä}{tehtävät}
\crefname{remark}{huomautus}{huomautukset}
\fi

\newcommand{\NN}{\mathbb{N}}
\newcommand{\NNone}{\mathbb{N}_1}
\newcommand{\ZZ}{\mathbb{Z}}
\newcommand{\QQ}{\mathbb{Q}}
\newcommand{\RR}{\mathbb{R}}
\newcommand{\CC}{\mathbb{C}}
\newcommand{\PP}{\mathbb{P}}

\DeclareMathOperator{\im}{im}
\DeclareMathOperator{\rank}{rank}
\DeclarePairedDelimiter{\abs}{\lvert}{\rvert}
\DeclarePairedDelimiter{\norm}{\lVert}{\rVert}
\newcommand{\set}[1]{\left\lbrace#1\right\rbrace}
\newcommand{\maxset}[1]{\max\set{#1}}
\newcommand{\minset}[1]{\min\set{#1}}
\newcommand{\powerset}[1]{\mathcal{P}\!\left(#1\right)}


\title{Analyysi I\\Ryhmätyö 5}
\author{}
\date{}

\begin{document}

\exerciseheading{Analyysi I}{Ryhmätyö 5}

\begin{exercise}
  \leavevmode\par
  \begin{enumerate}[label=(\alph*), nosep]
    \item Vertaa jonon ja funktion raja-arvon määritelmiä. Mitä eroa ja
      samankaltaisuuksia niissä on? Osaatko selittää matemaattisesti,
      mistä ne johtuvat?
    \item Anna esimerkki funktiosta \(f\colon \RR \to \RR\) ja kohti
      nollaa suppenevasta lukujonosta \((x_n)\), joille
      \(\lim_{x \to 0} f(x)\) ei ole olemassa mutta
      \(\lim_{n \to \infty} f(x_n)\) on olemassa (vrt. Lause 3.1.8).
  \end{enumerate}
\end{exercise}

\begin{solution}
  Olkoon \((x_n)\) lukujono, jolla \(x_n \to x\) ja \(f \colon A \to \RR\)
  funktio, jolle \(\lim_{t \to x_0} f(t) = L\).
  Merkitään pisteen \(a\in\RR\) \(\varepsilon\)-ympäristöä
  \[
    N_\varepsilon(a)=(a-\varepsilon,a+\varepsilon)
    =\{t\in\RR : |t-a|<\varepsilon\}, \qquad \varepsilon>0.
  \]
\begin{enumerate}[label=\alph*), nosep]
  \item Lukujonon tapauksessa \(x_n\to x\) tarkoittaa sitä, että jokaista
    \(\varepsilon>0\) kohti löytyy indeksi \(N\), jolle
    \(x_n\in N_\varepsilon(x)\) kaikilla \(n > N\).

    Funktion tapauksessa \(\lim_{t\to x_0}f(t)=L\) tarkoittaa sitä, että
    jokaista \(\varepsilon>0\) kohti löytyy \(\delta>0\) siten, että
    \[
      t\in (A\cap N_\delta(x_0))\setminus\{x_0\}
      \quad\Longrightarrow\quad f(t)\in N_\varepsilon(L).
    \]
    Molemmissa määritelmissä arvot saadaan siis mihin tahansa raja-arvon
    ympäristöön rajoittamalla syötteitä sopivasti. Lukujono on funktio,
    jonka määrittelyjoukko on \(\NN\): sen tapauksessa tarkastellaan
    riittävän suuria indeksejä. Funktion raja-arvossa pisteessä \(x_0\)
    taas tarkastellaan riittävän lähellä pistettä \(x_0\) olevia
    määrittelyjoukon pisteitä, mutta itse piste \(x_0\) jätetään pois.

    Funktion arvo pisteessä \(x_0\) ei vaikuta raja-arvoon, eikä funktion
    tarvitse edes olla määritelty siinä. Vastaavasti jonon äärellisen
    monen alkupään jäsenen muuttaminen ei vaikuta sen raja-arvoon.
  \item Määritellään \(f\colon\RR\to\RR\) asettamalla
    \[
      f(x)=
      \begin{cases}
        0, & x\le 0,\\
        1, & x>0.
      \end{cases}
    \]
    Valitaan \(x_n=1/n\), kun \(n\ge 1\). Tällöin \(x_n\to 0\) ja
    \(f(x_n)=1\) kaikilla \(n\), joten
    \(\lim_{n\to\infty}f(x_n)=1\).
    Toisaalta
    \[
      \lim_{x\to 0^-}f(x)=0
      \quad\text{ja}\quad
      \lim_{x\to 0^+}f(x)=1.
    \]
    Koska toispuoleiset raja-arvot ovat eri suuret,
    raja-arvoa \(\lim_{x\to 0}f(x)\) ei ole olemassa.
\end{enumerate}
\end{solution}

\begin{exercise}
Tutustukaa kirjassa Lauseen 3.1.8 todistukseen. Selittäkää todistuksen
idea noin kuudella ranskalaisella viivalle.
\end{exercise}

\begin{solution}
\leavevmode\par
\begin{itemize}[label=--, nosep]
  \item Todistuksessa osoitetaan funktion raja-arvon ja jonoehdon
    yhtäpitävyys, joten todistus sisältää kaksi suuntaa.
  \item Ensimmäisessä suunnassa oletetaan, että \(\lim_{x \to x_0} f(x) = L\),
    ja valitaan mielivaltainen funktion \(f\) määrittelyjoukon lukujono
    \((y_n)\), jolle \(y_n \to x_0\) ja \(y_n \neq x_0\) kaikilla \(n\).
    Tavoitteena on osoittaa, että \(f(y_n) \to L\).
  \item Kun \(\varepsilon>0\) on annettu, funktion raja-arvon määritelmä
    antaa sopivan luvun \(\delta>0\). Koska \(y_n \to x_0\), riittävän
    suurilla \(n\) pätee \(0<|y_n-x_0|<\delta\), joten
    \(|f(y_n)-L|<\varepsilon\). Siis \(f(y_n) \to L\).
  \item Toisessa suunnassa oletetaan, että \(f(y_n) \to L\) jokaisella
    määrittelyjoukon lukujonolla \((y_n)\), jolle \(y_n \to x_0\) ja
    \(y_n \neq x_0\) kaikilla \(n\). Tehdään vastaoletus, että
    \(f(x)\) ei suppene kohti lukua \(L\), kun \(x \to x_0\).
  \item Vastaoletuksen nojalla on olemassa \(\varepsilon_0>0\) siten,
    että jokaisella \(n\ge 1\) voidaan valita määrittelyjoukosta piste
    \(y_n\), jolle \(0<|y_n-x_0|<1/n\), mutta
    \(|f(y_n)-L|\ge\varepsilon_0\).
  \item Näin muodostettu jono toteuttaa ehdot \(y_n \to x_0\) ja
    \(y_n \neq x_0\), mutta \(f(y_n)\) ei suppene kohti lukua \(L\).
    Tämä on ristiriidassa oletuksen kanssa, joten vastaoletus on väärä
    ja \(\lim_{x \to x_0} f(x)=L\).
\end{itemize}
\end{solution}

\end{document}
